% \iffalse
%
% hamacro.dtx
% 俺のLaTeXマクロ集(煮ても焼いてもかまいません)
% hamano@cuspy.org
%
%<*driver>
\ProvidesFile{hamacro.dtx}
%</driver>
%<package>\NeedsTeXFormat{LaTeX2e}[1999/12/01]
%<package>\ProvidesPackage{hamacro}[2004/11/05 hamano]
%<*package>
%</package>
%
%<*driver>
%\documentclass{ltxdoc}
\documentclass[a4paper,12pt,uplatex]{jsarticle}
\usepackage{doc}
\usepackage[blue]{hamacro}
\usepackage{layout}
\EnableCrossrefs
\CodelineIndex
\RecordChanges
\setcounter{StandardModuleDepth}{1}

%\changes{v1.0}{2004/11/05}{Initial version}
\GetFileInfo{hamacro.dtx}
\DoNotIndex{\newcommand,\newenvironment}

\begin{document}
  \DocInput{hamacro.dtx}
  \PrintChanges
  \PrintIndex
\end{document}
%</driver>

% fontencを使うと何故かビットマップフォントが埋め込まれるので必ず欧文
% フォントを指定すること
\usepackage[T1]{fontenc}
\usepackage{newtxmath}
\renewcommand{\rmdefault}{ptm} % Nimbus Roman No9 L
\renewcommand{\sfdefault}{phv} % Nimbus Sans L
\renewcommand{\ttdefault}{zi4} % Inconsolata
\usepackage[deluxe]{otf}
\usepackage[dvipdfmx]{graphicx}

% 色定義
\usepackage[dvipdfmx]{color}
\usepackage[dvipdfmx]{xcolor}
\definecolor{lightgray}{HTML}{D3D3D3}
\definecolor{gainsboro}{HTML}{DCDCDC}
\definecolor{azure}{HTML}{F0FFFF}
\definecolor{darkblue}{HTML}{00008B}
\definecolor{midnightblue}{HTML}{191970}
\definecolor{lightskyblue}{HTML}{87CEFA}
\definecolor{olivegreen}{cmyk}{0.64,0,0.95,0.40}
\definecolor{cadetblue}{cmyk}{0.62,0.57,0.23,0}
\definecolor{brown}{cmyk}{0,0.81,1,0.60}

\definecolor{lightgreen}{HTML}{90EE90}
\definecolor{green}{HTML}{008000}
\definecolor{darkgreen}{HTML}{006400}
\definecolor{samuraiblue}{HTML}{1D327D}

\usepackage[dvipdfm]{hyperref}
\hypersetup{
  setpagesize=false,
  bookmarksnumbered=true,
  colorlinks=true,
  linkcolor=linkcolor,
  citecolor=linkcolor,
  filecolor=linkcolor,
  urlcolor=linkcolor
}

% しおりの文字化け対策
\usepackage{pxjahyper}

% その他パッケージ
\usepackage{here}
\usepackage{ctable}
\usepackage{longtable}
\usepackage{ulem}
\usepackage{soul}

% カラーテーマ設定
\colorlet{framerule}{gray}
\colorlet{framebg}{gainsboro}
\colorlet{linkcolor}{blue}

\colorlet{accentcolor}{black}
\colorlet{subcolor}{black}

\DeclareOption{green}{
  \colorlet{accentcolor}{green}
  \colorlet{subcolor}{darkgreen}
}
\DeclareOption{blue}{
  \colorlet{accentcolor}{lightskyblue}
  \colorlet{subcolor}{midnightblue}
}
\DeclareOption{samuraiblue}{
  \colorlet{accentcolor}{lightskyblue}
  \colorlet{subcolor}{samuraiblue}
}
\ProcessOptions

% jsclassesとgeometryパッケージは相性が悪い
%\usepackage[a4paper,top=35mm, bottom=30mm, left=30mm, right=30mm]{geometry}
\setlength{\oddsidemargin}{-1truein}
\addtolength{\oddsidemargin}{0.11\paperwidth}
\setlength{\evensidemargin}{-1truein}
\addtolength{\evensidemargin}{0.11\paperwidth}
\setlength{\marginparwidth}{0.10\paperwidth}
\setlength{\marginparsep}{0.005\paperwidth}
\setlength{\textwidth}{\paperwidth}
\addtolength{\textwidth}{-0.22\paperwidth}

\message{^^J}
\message{***** Layout Info *****^^J}
\message{\hoffset: \the\hoffset^^J}
\message{\voffset: \the\voffset^^J}
\message{\paperwidth: \the\paperwidth^^J}
\message{\paperheight: \the\paperheight^^J}
\message{\textwidth: \the\textwidth^^J}
\message{\textheight: \the\textheight^^J}
\message{\oddsidemargin: \the\oddsidemargin^^J}
\message{\evensidemargin: \the\evensidemargin^^J}
\message{\marginparwidth: \the\marginparwidth^^J}
\message{\marginparsep: \the\marginparsep^^J}
\message{***********************^^J}

% default title formatting
\def\@maketitle{
  \newpage
  \null
  \begin{center}
    {\LARGE\gtfamily\bfseries \@title \par}
  \end{center}
  \begin{flushright}
    {\large \@date \par}
    {\large \@author \par}
  \end{flushright}
  \vspace{0.5in}
}

% chapter formatting
\usepackage[Conny]{fncychap}
\ChNameVar{\centering\Huge\sffamily\color{black}}
\ChTitleAsIs
\ChTitleVar{\centering\Huge\sffamily\color{black}}

\renewcommand{\DOCH}{
  \color{subcolor}
  \mghrulefill{3\RW}\par\nobreak
  \vskip -0.5\baselineskip
  \mghrulefill{\RW}\par\nobreak
  \CNV\FmN{\@chapapp}\space \CNoV\thechapter \postchaptername
  \par\nobreak
  \vskip -0.5\baselineskip
}
\renewcommand{\DOTI}[1]{
  \color{subcolor}
  \mghrulefill{\RW}\par\nobreak
  \CTV\FmTi{#1}\par\nobreak
  \vskip 1\baselineskip
}
\renewcommand{\DOTIS}[1]{
  \color{subcolor}
  \mghrulefill{\RW}\par\nobreak
  \CTV\FmTi{#1}\par\nobreak
  \vskip 1\baselineskip
}

% section formatting
\usepackage[explicit]{titlesec}
\titleformat{name=\section}[block]
{\gtfamily\bfseries\LARGE}
{\color{subcolor}{\thesection} \color{black}{#1}}
{1em}{}
[\titleline{\color{subcolor}\titlerule[1.5pt]}]

\titleformat{name=\section,numberless}[block]
{\gtfamily\bfseries\LARGE}
{#1}
{1em}{}
[\titleline{\color{subcolor}\titlerule[1.5pt]}]

% subsection formatting
\newcommand\subseculine{\bgroup\markoverwith
  {\textcolor{subcolor}{\rule[-1ex]{2pt}{1.5pt}}}\ULon}

\titleformat{\subsection}[block]
{\gtfamily\bfseries\Large}
{\subseculine{\textcolor{subcolor}{\thesubsection{}} {#1}}}
{1em}{}

\titleformat{name=\subsection,numberless}[block]
{\gtfamily\bfseries\Large}
{\subseculine{#1}}
{1em}{}

% subsubsection formatting
\titleformat{\subsubsection}[block]
{\gtfamily\bfseries\large}
{\color{subcolor}{\thesubsubsection} \color{black}{#1}}{10em}{}

\titleformat{name=\subsubsection,numberless}[block]
{\gtfamily\bfseries\large}
{\color{black}{#1}}{10em}{}

% 修飾マクロ
\let\origemph\emph
\newcommand{\redemph}[1]
           {\textcolor{red}{\underline{\origemph{\textcolor{black}{#1}}}}}

\newcommand{\ulemph}[1]
           {\textbf{\Underline{#1}}}

\newcommand{\bouemph}[1]
           {\origemph{\bou{#1}}}

%\let\emph\bouemph

% style for quote
\usepackage{framed}
\usepackage{anyfontsize}
\newcommand{\bigquotemark}{
  {\fontsize{32pt}{32pt}\selectfont \color{accentcolor} \textbf{``}}
  \vskip -2.3em
  \hskip 0.5zw
}
\renewenvironment{quote}{
  \def\FrameCommand{
    \fboxsep=\FrameSep
    \fcolorbox{gray}{gainsboro}
  }
  \MakeFramed {\advance\hsize-\width \FrameRestore}
  \parindent 0zw
  \vskip -0.5em
  \hskip -1zw
  {\bigquotemark}
}
{
  \smallskip
  \endMakeFramed
}

% style for verbatim
\usepackage{fancyvrb}

\renewenvironment{verbatim}{
  \VerbatimEnvironment
  \def\FrameCommand{
    \fboxsep=4truept
    \fcolorbox{gray}{gainsboro}
  }
  \MakeFramed{\advance\hsize-\width \FrameRestore}
  \tt\fontsize{11.2truept}{1.2em}
  \begin{Verbatim}}
{
  \end{Verbatim}
  \endMakeFramed
}

\newenvironment{verbatim60}{
  \VerbatimEnvironment
  \def\FrameCommand{
    \fboxsep=4truept
    \fcolorbox{gray}{gainsboro}
  }
  \MakeFramed{\advance\hsize-\width \FrameRestore}
  \tt\fontsize{15truept}{1.5em}
  \begin{Verbatim}}
{
  \end{Verbatim}
  \endMakeFramed
}

\newenvironment{verbatim40}{
  \VerbatimEnvironment
  \def\FrameCommand{
    \fboxsep=4truept
    \fcolorbox{gray}{gainsboro}
  }
  \MakeFramed{\advance\hsize-\width \FrameRestore}
  \tt\fontsize{22truept}{1.5em}
  \begin{Verbatim}}
{
  \end{Verbatim}
  \endMakeFramed
}

% style for pandoc block
\usepackage{mdframed}
\newenvironment{Shaded}{\begin{mdframed}}{\end{mdframed}}
\DefineVerbatimEnvironment{Highlighting}{Verbatim}{commandchars=\\\{\}}
\newcommand{\KeywordTok}[1]{\textcolor[rgb]{0.00,0.44,0.13}{\textbf{{#1}}}}
\newcommand{\DataTypeTok}[1]{\textcolor[rgb]{0.56,0.13,0.00}{{#1}}}
\newcommand{\DecValTok}[1]{\textcolor[rgb]{0.25,0.63,0.44}{{#1}}}
\newcommand{\BaseNTok}[1]{\textcolor[rgb]{0.25,0.63,0.44}{{#1}}}
\newcommand{\FloatTok}[1]{\textcolor[rgb]{0.25,0.63,0.44}{{#1}}}
\newcommand{\CharTok}[1]{\textcolor[rgb]{0.25,0.44,0.63}{{#1}}}
\newcommand{\StringTok}[1]{\textcolor[rgb]{0.25,0.44,0.63}{{#1}}}
\newcommand{\CommentTok}[1]{\textcolor[rgb]{0.38,0.63,0.69}{\textit{{#1}}}}
\newcommand{\OtherTok}[1]{\textcolor[rgb]{0.00,0.44,0.13}{{#1}}}
\newcommand{\AlertTok}[1]{\textcolor[rgb]{1.00,0.00,0.00}{\textbf{{#1}}}}
\newcommand{\FunctionTok}[1]{\textcolor[rgb]{0.02,0.16,0.49}{{#1}}}
\newcommand{\RegionMarkerTok}[1]{{#1}}
\newcommand{\ErrorTok}[1]{\textcolor[rgb]{1.00,0.00,0.00}{\textbf{{#1}}}}
\newcommand{\NormalTok}[1]{{#1}}

% fix for pandoc 1.14
\providecommand{\tightlist}{
  \setlength{\itemsep}{0pt}\setlength{\parskip}{0pt}}

% style for listings
\usepackage{upquote}
\usepackage{listings}
%\usepackage{jlisting}

% ハイフンをU+002D HYPHEN-MINUSのままで表示
%\makeatletter
%\lst@CCPutMacro\lst@ProcessOther {"2D}{\lst@ttfamily{-{}}{-{}}}
%\@empty\z@\@empty
%\makeatother

% アスタリスクをU+002A ASTERISKのままで表示
%\makeatletter
%\lst@CCPutMacro\lst@ProcessOther {"2A}{\lst@ttfamily{*{}}{*{}}}
%\@empty\z@\@empty
%\makeatother

% 一部のPDFリーダーで白い線が入ってしまう問題があるのでmdframeでラップする
% http://tex.stackexchange.com/questions/196247/lstlisting-width-of-background-box

% captionもmdframe内に入るのでいまいち
\AtBeginDocument{
  \def\lstlistingname{リスト}
}

\mdfsetup{
  linewidth=1truept,
  linecolor=gray,
  innerleftmargin=0.5zw,
  innerrightmargin=0.5zw,
  backgroundcolor=gainsboro
}

\BeforeBeginEnvironment{lstlisting}{
  \begin{mdframed}
}

\AfterEndEnvironment{lstlisting}{
  \end{mdframed}
}

\lstset{
  frame=single,
  rulecolor=\color{gainsboro},
  backgroundcolor=\color{gainsboro},
%  numbers=left,
%  numbersep=0pt,
  framesep=0pt,
  xleftmargin=0pt,
  basicstyle=\fontsize{11.2truept}{1.2em}\tt,
  belowskip=-0.5em,
  identifierstyle=\color{black},
  commentstyle=\color{brown},
  keywordstyle=\color{olivegreen},
  ndkeywordstyle=\color{red},
  stringstyle=\color{purple},
  showstringspaces=false,
  breaklines=true,
  columns=[l]{fullflexible}
}

\lstdefinestyle{command}{
  language=command,
  commentstyle=\color{black}
}

\lstdefinelanguage{command}
{
  keywords={export, if, while, for},
  morekeywords={\$, \#, \%},
  morestring=[b]',
  morestring=[b]\"
}

\lstdefinelanguage{JavaScript}{
  keywords={typeof, new, true, false, catch, function, return, null, catch, switch, var, if, in, while, do, else, case, break},
  ndkeywords={class, export, boolean, throw, implements, import, this},
  sensitive=false,
  comment=[l]{//},
  morecomment=[s]{/*}{*/},
  morestring=[b]',
  morestring=[b]\"
}

\lstdefinelanguage{httpd.conf}
{
  morekeywords={VirtualHost, Include, Directory, LogFormat, CustomLog},
  morecomment=[l]{\#}
}

\lstdefinelanguage{nginx.conf}
{
  morekeywords={
    server,server_name,listen,location,
    rewrite,proxy_pass, proxy_set_header},
  morecomment=[l]{\#}
}

% \fi
%
% \iffalse
%<*doc>
% \fi

\title{\LaTeX{} hamacro パッケージ}
\author{hamano}
\maketitle
\tableofcontents
\clearpage
%\MakeShortVerb{\|}

\section{はじめに}

これは私が本や文書を組版する際に利用する\LaTeX{}のスタイルパッケージで
す。

80文字ピッタリおさまるverbatim環境や、各種プログラミング言語のコードハ
イライティングなど、主にプログラミングに関する文書に適しています。

既存の\LaTeX{}文書で簡単に利用できることと、pandocやRe:Viewが出力する
\LaTeX{}で利用することを想定していますので、新しいマクロを定義するとい
うよりも既存のマクロや環境を書き換えています。

このパッケージはMITライセンスで配布していますので、利用、改変、部分的な
コピーなど自由です!

\section{使い方}

hamacro.styを然るべき場所に配置し、プリアンブルに

\begin{lstlisting}[language=tex]
\usepackage{hamacro}
\end{lstlisting}

と記述するだけでこのドキュメントの様なスタイルで組版できます。

具体的な用例は\href{https://github.com/hamano/hamacro/blob/master/example.tex}{example.tex}を参照して下さい。

\subsection{pandocでの利用方法}

このスタイルはpandocで生成したTeX文書に最適化しています。
ですのでmarkdownで書いた文書を手軽にPDFに組版できます。
例えば\href{https://github.com/hamano/hamacro/blob/master/examples/memo/memo.pdf}{memo.md}を変換して\href{https://github.com/hamano/hamacro/blob/master/examples/memo/memo.pdf}{memo.pdf}を生成できます。

詳しくはこちらのサンプルディレクトリを見てください。

 - \href{https://github.com/hamano/hamacro/blob/master/examples/memo/}{https://github.com/hamano/hamacro/blob/master/examples/memo/}

\subsection{Re:Viewでの利用方法}

\section{用例}
\subsection{フォント}

\subsubsection{deluxeフォント}
\begin{itemize}
  \item {\mcfamily 明朝体(ヒラギノ明朝W3)}
  \item {\mcfamily\bfseries 明朝太字(ヒラギノ明朝W6)}
  \item {\gtfamily ゴシック(ヒラギノ角ゴW3)}
  \item {\gtfamily\bfseries ゴシック太字(ヒラギノ角ゴW6)}
  \item {\gtfamily\ebseries ゴシック極太(ヒラギノ角ゴW8)}
  \item {\mgfamily 丸文字(ヒラギノ丸ゴW4)}
  \item {\mcfamily\ltseries 明朝細字(ヒラギノ明朝W2)}
\end{itemize}

\subsubsection{書体}
\begin{itemize}
  \item ファミリー
  \begin{itemize}
    \item \textrm{ローマン体}
    \item \textsf{サンセリフ体}
    \item \texttt{タイプライター体}
    \item \textmc{明朝体}
    \item \textgt{ゴシック}
  \end{itemize}

  \item シリーズ
  \begin{itemize}
    \item \textmd{ミディアム}
    \item \textbf{ボールド}
  \end{itemize}
\end{itemize}

\subsubsection{文字修飾}
\begin{itemize}
\item 普通の文字(normal)
\item \emph{強調(emph)}
\item \sout{打ち消し線}
\item \textsuperscript{上付き}
\end{itemize}

\subsubsection{フォントサイズ}
\def\sampletext{Windowsでコンピュータの世界が広がります。}
\begin{itemize}
  \normalsize\item\tiny{\sampletext(tiny)}
  \normalsize\item\scriptsize{\sampletext(scriptsize)}
  \normalsize\item\footnotesize{\sampletext(footnotesize)}
  \normalsize\item\small{\sampletext(small)}
  \normalsize\item\normalsize{\sampletext(normalsize)}
  \normalsize\item\large{\sampletext(large)}
  \normalsize\item\Large{\sampletext(Large)}
  \normalsize\item\LARGE{\sampletext(LARGE)}
  \normalsize\item\huge{\sampletext(huge)}
  \normalsize\item\Huge{\sampletext(Huge)}
\end{itemize}
\normalsize

\subsection{リンク}
\url{http://localhost/}

\href{http://www.example.com/}{example.com}

\subsection{quote環境}

\begin{quote}
吾輩は猫である。名前はまだ無い。

どこで生れたかとんと見当がつかぬ。
\end{quote}

\subsection{quotation環境}

\begin{quotation}
吾輩は猫である。名前はまだ無い。

どこで生れたかとんと見当がつかぬ。
\end{quotation}

\subsection{verbatim環境}

半角80文字がちょうど収まるサイズにしています。
コマンドラインやソースコードを示す場合はverbatim環境よりもlistingsを推奨します

\begin{verbatim}
01234567890123456789012345678901234567890123456789012345678901234567890123456789
０１２３４５６７８９０１２３４５６７８９０１２３４５６７８９０１２３４５６７８９
\end{verbatim}

フォントが小さくなって嫌だ、という方のために60文字がちょうど収まるverbatim60とverbatim40を用意しています。

\subsection{verbatim60}

\begin{verbatim60}
012345678901234567890123456789012345678901234567890123456789
０１２３４５６７８９０１２３４５６７８９０１２３４５６７８９
\end{verbatim60}

\subsection{verbatim40}

\begin{verbatim40}
0123456789012345678901234567890123456789
０１２３４５６７８９０１２３４５６７８９
\end{verbatim40}

\subsection{コードリスティング}

\subsubsection{コマンドライン}
\begin{lstlisting}[language=command]
# uname --machine
x86_64
\end{lstlisting}

\subsubsection{C}
\begin{lstlisting}[language=C]
#include <stdio.h>

int main(int argc, char *argv[]){
    printf("Hello World!");
    return 0;
}
\end{lstlisting}

\subsubsection{Python}
\begin{lstlisting}[language=python]
def hello():
    print('Hello World!')

if __name__ == '__main__':
    hello()
\end{lstlisting}

\subsubsection{80文字テスト}
\begin{lstlisting}
01234567890123456789012345678901234567890123456789012345678901234567890123456789
\end{lstlisting}

\subsection{リスト}

\subsubsection{列挙リスト}
\begin{itemize}
\item りんご
\item みかん
\item バナナ
\end{itemize}

\subsubsection{順序リスト}
\begin{enumerate}
\item 睦月
\item 如月
\item 弥生
\end{enumerate}

\subsubsection{定義リスト}
\begin{description}
\item[Monday] 月曜日
\item[Tuesday] 火曜日
\item[Wednesday] 水曜日
\end{description}

\subsection{テーブル}

\subsubsection{普通のテーブル}
\begin{table}[htbp]
\begin{center}
\begin{tabular}{|l|c|r|}
\hline
セル1 & セル2 & セル3 \\ \hline
セル4 & セル5 & セル6 \\ \hline
セル7 & セル8 & セル9 \\ \hline
\end{tabular}
\caption{キャプション}
\end{center}
\end{table}

\subsubsection{longtable}
\begin{longtable}[c]{@{}crr@{}}
\toprule\addlinespace
果物 & 値段 & 個数
\\\addlinespace
\midrule\endhead
バナナ & 134円 & 10本
\\\addlinespace
オレンジ & 210円 & 10個
\\\addlinespace
\bottomrule
\end{longtable}

\subsection{水平線}
\begin{center}\rule{3in}{0.4pt}\end{center}

\subsection{画像}
\includegraphics[width=\textwidth]{latex.eps}

\clearpage
\subsection{数式}
\begin{enumerate}

\item 総和
\begin{equation}
\sum_{i=1}^n i = \frac{n(n+1)}{2}
\end{equation}

\item 質量とエネルギーの等価性
\begin{equation}
E = mc ^2
\end{equation}

\item 波動方程式
\begin{equation}
\frac{\partial^2 z}{\partial t^2}=
c^2 (\frac{\partial^2 z}{\partial x^2}+\frac{\partial^2 z}{\partial y^2})-
\mu \frac{\partial z}{\partial t}
\end{equation}

\item 文章中に突然の $e^{i\theta}=\cos\theta+i\sin\theta$ オイラーの公式

\end{enumerate}


\section{レイアウト}
\layout

% \iffalse
%</doc>
% \fi
%
% \Finale
\endinput
%
% \CheckSum{0}
%
% \CharacterTable
%  {Upper-case    \A\B\C\D\E\F\G\H\I\J\K\L\M\N\O\P\Q\R\S\T\U\V\W\X\Y\Z
%   Lower-case    \a\b\c\d\e\f\g\h\i\j\k\l\m\n\o\p\q\r\s\t\u\v\w\x\y\z
%   Digits        \0\1\2\3\4\5\6\7\8\9
%   Exclamation   \!     Double quote  \"     Hash (number) \#
%   Dollar        \$     Percent       \%     Ampersand     \&
%   Acute accent  \'     Left paren    \(     Right paren   \)
%   Asterisk      \*     Plus          \+     Comma         \,
%   Minus         \-     Point         \.     Solidus       \/
%   Colon         \:     Semicolon     \;     Less than     \<
%   Equals        \=     Greater than  \>     Question mark \?
%   Commercial at \@     Left bracket  \[     Backslash     \\
%   Right bracket \]     Circumflex    \^     Underscore    \_
%   Grave accent  \`     Left brace    \{     Vertical bar  \|
%   Right brace   \}     Tilde         \~}
%
%
